% ==STYLE-BEGIN==
\documentclass[UTF8,11pt]{ctexart}
\usepackage{amsmath,amssymb,amsthm}
\usepackage[a4paper,margin=2.4cm,top=2.3cm,bottom=2.3cm]{geometry}
\usepackage{fancyhdr}
\usepackage{enumitem}
\usepackage{hyperref}
\hypersetup{hidelinks}

% ---- 版式 ----
\linespread{1.15}
\setlength{\parskip}{0.35em}
\setlength{\parindent}{1.6em}
\setlist{nosep,leftmargin=2em}

% ---- 页眉 ----
\pagestyle{fancy}
\fancyhf{}
\fancyhead[L]{\small\coursename}
\fancyhead[R]{\small\thepage}
\renewcommand{\headrulewidth}{0.4pt}

% ---- 定理环境：中文用正体，标题加粗，按小节共用一条流水号 ----
\newtheoremstyle{cnthm}{0.7em}{0.7em}{\normalfont}{}{\bfseries}{}{0.5em}{}
\theoremstyle{cnthm}
\newtheorem{thm}{定理}[section]
\newtheorem{lem}[thm]{引理}
\newtheorem{prop}[thm]{命题}
\newtheorem{cor}[thm]{推论}
\newtheorem{defn}[thm]{定义}
\newtheorem{rem}[thm]{注}
\newtheorem{exa}[thm]{例}
\newtheorem{conj}[thm]{猜想}
\renewcommand{\proofname}{\bfseries 证明}
\providecommand{\coursename}{}
% ==STYLE-END==
\renewcommand{\coursename}{解析几何\,·\,0924}
\title{坐标变换公式}
\xeCJKDeclareCharClass{CJK}{"2150 -> "218F} % 罗马数字 Ⅰ Ⅱ Ⅲ 使用中文字体
\date{}
\begin{document}
\maketitle

\section{\ 设定}

$[O;\vec e_1,\vec e_2,\vec e_3]$ 为直角坐标系（Ⅰ），
$[O';\vec e_1',\vec e_2',\vec e_3']$ 为直角坐标系（Ⅱ）。

$O'$ 在旧坐标系中的坐标是 $(a_1,a_2,a_3)$；

$\vec e_1'$ 在旧坐标系中的坐标是 $(a_{11},a_{21},a_{31})$，
$\vec e_2'$ 在旧坐标系中的坐标是 $(a_{12},a_{22},a_{32})$，
$\vec e_3'$ 在旧坐标系中的坐标是 $(a_{13},a_{23},a_{33})$，即
\[
\begin{cases}
\vec e_1'=a_{11}\vec e_1+a_{21}\vec e_2+a_{31}\vec e_3,\\
\vec e_2'=a_{12}\vec e_1+a_{22}\vec e_2+a_{32}\vec e_3,\\
\vec e_3'=a_{13}\vec e_1+a_{23}\vec e_2+a_{33}\vec e_3,
\end{cases}
\qquad
\overrightarrow{OO'}=a_1\vec e_1+a_2\vec e_2+a_3\vec e_3 .
\]

此时，设点 $P$ 在Ⅰ、Ⅱ的坐标分别为 $(x,y,z)^T$ 和
$(x',y',z')^T$。试问：$(x,y,z)^T$ 和 $(x',y',z')^T$
关系如何？

\section{\ 用新坐标表示旧坐标}

易得
\[
\overrightarrow{OP}=x\vec e_1+y\vec e_2+z\vec e_3,\qquad
\overrightarrow{O_1P_1}
=x'\vec e_1'+y'\vec e_2'+z'\vec e_3' .
\]
我们将 $\vec e_1',\vec e_2',\vec e_3'$ 全部替换为
$\vec e_1,\vec e_2,\vec e_3$，可以得到下式
\[
\begin{aligned}
\overrightarrow{O_1P_1}
&=\ x'(a_{11}\vec e_1+a_{21}\vec e_2+a_{31}\vec e_3)\\
&\quad+\ y'(a_{12}\vec e_1+a_{22}\vec e_2+a_{32}\vec e_3)\\
&\quad+\ z'(a_{13}\vec e_1+a_{23}\vec e_2+a_{33}\vec e_3)\\
&\quad+\ \overrightarrow{OO_1},
\end{aligned}
\]
整理可得
\[
\begin{aligned}
&=\ (x'a_{11}+y'a_{12}+z'a_{13})\vec e_1\\
&\quad+\ (x'a_{21}+y'a_{22}+z'a_{23})\vec e_2\\
&\quad+\ (x'a_{31}+y'a_{32}+z'a_{33})\vec e_3\\
&\quad+\ \overrightarrow{OO_1} .
\end{aligned}
\]
不难发现
\[
\begin{cases}
x=x'a_{11}+y'a_{12}+z'a_{13}+a_1,\\
y=x'a_{21}+y'a_{22}+z'a_{23}+a_2,\\
z=x'a_{31}+y'a_{32}+z'a_{33}+a_3 .
\end{cases}
\]
在这里，我们引入矩阵与向量的乘法，可以将上式写成如下形式
\[
\begin{pmatrix}
x\\y\\z
\end{pmatrix}
=
\begin{pmatrix}
a_{11} & a_{12} & a_{13}\\
a_{21} & a_{22} & a_{23}\\
a_{31} & a_{32} & a_{33}
\end{pmatrix}
\begin{pmatrix}
x'\\y'\\z'
\end{pmatrix}
+
\begin{pmatrix}
a_1\\a_2\\a_3
\end{pmatrix}.
\]
我们记中间的矩阵为 $A$，则上式可写为
\[
X=AX'+X_0 . \tag{1}
\]

\section{\ 用旧坐标表示新坐标}

以上，我们达成了用新坐标表示旧坐标的目的，但实际算题的时候，
我们更希望通过旧坐标去表示新坐标，那又该怎么办呢？

回忆投影的概念，坐标实际上就是 $\overrightarrow{OP}$ 在
$\vec e_1,\vec e_2,\vec e_3$ 上的投影，我们可以就此思路去表示
$x',y',z'$，下以 $x'$ 为例
\[
\begin{aligned}
x'&=\overrightarrow{O_1P}\cdot\vec e_1'
=(\overrightarrow{O_1O}+\overrightarrow{OP})\cdot\vec e_1'\\
&=(x-a_1,\ y-a_2,\ z-a_3)(a_{11},\ a_{21},\ a_{31}),
\end{aligned}
\]
有
\[
x'=a_{11}(x-a_1)+a_{21}(y-a_2)+a_{31}(z-a_3).
\]
可以得到
\[
\begin{pmatrix}
x'\\y'\\z'
\end{pmatrix}
=
\begin{pmatrix}
a_{11} & a_{21} & a_{31}\\
a_{12} & a_{22} & a_{32}\\
a_{13} & a_{23} & a_{33}
\end{pmatrix}
\left(
\begin{pmatrix}
x\\y\\z
\end{pmatrix}
-
\begin{pmatrix}
a_1\\a_2\\a_3
\end{pmatrix}
\right).
\]
这个矩阵是上面矩阵的转置，记作 $A^T$，可以写为下式
\[
X'=A^T(X-X_0). \tag{2}
\]

这个式子和 (1) 式高度相似，矩阵 $A$ 好似被“除”到了另一面。
矩阵没有除法，只有乘法，为达成“除”的目的，
我们引入矩阵逆阵的概念
\[
AA^{-1}=I .
\]
这样，我们将 (1) 式两边同时乘 $A^{-1}$ 再移项，
就能化成类似 (2) 式的形式如下
\[
X'=A^{-1}(X-X_0).
\]
两种方法得到的式子理应是一样的，即 $A^T=A^{-1}$。
我们称这样的矩阵是“\textbf{正交矩阵}”。

\section{\ 连续变换与矩阵乘法}

如果此时有第三个坐标系，我们要从Ⅰ系连续变换到Ⅱ系再到Ⅲ系呢？

我们暂时不考虑坐标系平移，只考虑旋转，对于Ⅰ变换到Ⅱ、
Ⅱ变换到Ⅲ我们分别有
\[
(\mathrm{I})\to(\mathrm{II}):\quad X=A\cdot X',\qquad A=(a_{ij});
\]
\[
(\mathrm{II})\to(\mathrm{III}):\quad X'=B\cdot X'',\qquad B=(b_{ij}).
\]
将下式代入上式即得
\[
(\mathrm{I})\to(\mathrm{III}):\quad X=AB\,X'' .
\]
我们引入了矩阵之间的乘法。

\section{\ 正交基底带来的性质}

回忆我们一开始的假设，我们假设
$\vec e_1,\vec e_2,\vec e_3$ 是一组正交基底，
这会导致这个矩阵有很多奇妙的性质。

由于 $\vec e_i\cdot\vec e_j=\delta_{ij}$，
$\forall i,j=1,2,3$，有
\[
\sum_{k=1}^{3}a_{ki}a_{kj}=\delta_{ij},\qquad
\sum_{k=1}^{3}a_{ik}a_{jk}=\delta_{ij},
\qquad\forall i,j=1,2,3 .
\]
事实上，上面两式等价于矩阵乘法
\[
A^TA=I,\qquad AA^T=I .
\]

\end{document}
